\section{相关工作}

\subsection{区域风电功率预测与时空建模}

风电功率预测可以按照空间尺度划分为单场站预测、场群预测和区域预测。对于电力系统调度而言，目标区域内多个风电场站的聚合出力通常比单个场站出力更直接对应功率平衡需求。区域或场群尺度的预测不仅需要描述功率序列的时间变化，还需要利用不同场站之间的空间相关性以及区域天气过程的连续传播特征。

已有研究从多种时空建模角度提升了风电预测精度。早期深度学习方法利用 CNN-LSTM 等结构联合建模局部空间特征和时间依赖 \cite{ref3}，也有工作将风机或场站嵌入规则网格，使卷积网络能够处理空间分布信息 \cite{ref4}。近年来，图神经网络被广泛用于场群级风电预测，因为风电场站可以自然表示为图节点，节点之间的边可由地理距离、气象相关性或功率相关性构造。相关方法包括时空图神经网络与模式识别结合 \cite{ref5}、全局信息自适应感知图卷积 \cite{ref6}、大规模场群动态图建模 \cite{ref7}、可学习图结构和动态误差修正 \cite{ref8}、趋势切换机制 \cite{ref9}、双融合时空图模型 \cite{ref10}、趋势分类与空间信息集成 \cite{ref11}，以及区域短期预测中的层级图神经网络 \cite{ref2} 和持续学习时空图框架 \cite{ref12}。

这些研究证明了多场站时空相关性建模对区域或场群预测的重要性。然而，现有方法的重点通常放在图结构构造、历史功率动态或场站间空间依赖上；即使引入气象信息，也常被作为节点特征或辅助输入使用。对于业务化多步预测中未来气象预报和近期场站功率的不同角色，即未来气象作为预测窗口内主要外生驱动、近期功率作为目标窗口前状态条件，已有工作讨论仍不充分。

\subsection{气象预报驱动的风电功率预测}

在日内多步风电预测中，未来气象预报是不可替代的外生输入。传统数值天气预报（Numerical Weather Prediction, NWP）和深度气象预报模型虽然来源不同，但都可以提供预测窗口内的风场和背景气象信息。本文使用 Pangu-Weather \cite{ref22} 生成的网格化气象预报作为输入，其建模作用与 NWP-功率融合研究中的未来气象驱动一致。

已有方法通常将 NWP 与历史功率作为两类输入，通过特征拼接、并行编码器、多模态注意力或多任务学习进行融合。例如，M2TNet \cite{ref13} 和 M2STAN \cite{ref14} 将 NWP 与功率信息建模为多模态多任务信息流；围绕 NWP 本身的误差特征，已有研究讨论了 NWP 预报风速与实测风速之间的时间滞后 \cite{ref15}、NWP 风速误差容忍性 \cite{ref16}、多源时空特征融合 \cite{ref17}，以及校正 NWP 与空间相关性建模的结合 \cite{ref18}。除节点或序列形式的 NWP 特征外，也有研究直接编码网格化气象场，例如使用 3D CNN 提取 NWP 数据立方体特征 \cite{ref19}、采用自动化深度学习框架进行区域风电预测 \cite{ref20}，或将深度 CNN 与概率预测框架结合用于区域概率风电预测 \cite{ref21}。近期工作还探索了大语言模型在多模态风电预测中的应用 \cite{ref23}。

上述研究表明，未来气象预报对区域风电预测具有关键作用，也说明直接编码网格化气象场是有效的建模方向。但多数融合框架仍将气象信息和功率信息作为相对平行的输入分支，通过拼接、并行编码或注意力交互进行对称融合。本文关注的是另一种角色划分：未来网格化气象预报作为预测窗口内的主要驱动，近期场站功率作为当前状态条件，对未来气象记忆的读取过程进行约束。

\subsection{条件化记忆与非对称跨模态融合}

Transformer 及其变体为多步时间序列预测提供了重要基础。Informer \cite{ref24}、PatchTST \cite{ref27}、Crossformer \cite{ref28} 和 Non-stationary Transformers \cite{ref34} 分别从长序列注意力、时间 patch 表示、跨变量依赖和非平稳性角度改进序列预测。风电预测中也已有多种 Transformer 或图 Transformer 结构用于建模时空耦合关系 \cite{ref37,ref38,ref39,ref40,ref41,ref42,ref43}。这些方法说明，注意力机制和并行多步解码适合处理复杂的风电时序预测问题。

与本文解码器设计更直接相关的是 query-based decoding 和可被注意力读取的记忆结构。DETR 使用 learnable queries 通过 cross-attention 读取编码特征并生成结构化输出 \cite{ref35}；Perceiver 和 Perceiver IO 通过 latent array 与 query-based decoding 处理大规模结构化输入输出 \cite{ref65,ref36}。更早的 Memory Networks 和 End-to-End Memory Networks 也说明，可读写或可注意力访问的记忆表示能够作为模型推理的信息载体 \cite{ref60,ref61}。在多模态系统中，Flamingo 使用视觉 tokens 注入语言模型，FiLM 则通过条件化调制改变特征表示 \cite{ref66,ref70}。这些工作为“用条件信息影响主模态读取过程”提供了形式参照。

需要强调的是，本文并不进行 prompt learning、语言模型适配或外部检索增强。本文中的 dynamic power context tokens 由当前样本的场站级功率动态直接生成，并被注入气象记忆空间；future queries 读取的是由功率状态条件化后的气象记忆。与风电预测中常见的特征级注意力或双注意力结构 \cite{ref72,ref73} 不同，本文的关键在于将近期功率作为状态条件前置到气象记忆中，使预测查询在读取未来气象驱动时显式受到当前多场站运行状态的约束。因此，本文方法更准确地表述为面向区域风电预测的功率条件化气象记忆，而不是通用 prompt 方法或对称多模态融合。
